
\chapter{Subregular and hook-type \texorpdfstring{$W$}{W}-algebras from first principles}
\label{chap:subregular-hook-frontier}

\section{The governing question}

The subregular reduction of $V^k(\mathfrak{sl}_N)$ is the first
non-principal family for which one can place the quantum
Drinfeld--Sokolov presentation, the strong-generator OPE, and the
candidate bar companion on the same page.  These three objects carry
different kinds of information.  The Kazhdan--Li filtration controls
the associated graded vertex Poisson algebra; the Feigin--Semikhatov
realisation controls the singular OPE coefficients; the completed
chiral bar construction controls twisting morphisms and derived
reconstruction.  A theorem connecting them includes the comparison
map as part of its hypotheses.

The chapter therefore has two mathematical surfaces.

\begin{enumerate}[label=(\roman*)]
\item The \emph{exact OPE surface} consists of the filtered Slodowy
 presentation, the Feigin--Semikhatov singular products, the Bell recursion,
 and the centred Miura--Appell symbols.  These statements are algebraic and
 are proved below or cited from the primary vertex-algebra literature.
\item The \emph{derived comparison surface} asks whether those filtered data
 determine the completed chiral bar homology, its Verdier dual, Theorem~H,
 and a same-family companion.  Each such assertion is conditional on a named
 convergence, comparison, and nondegeneracy package.
\end{enumerate}

The load-bearing distinction is visible in the following table.
\begin{center}
\renewcommand{\arraystretch}{1.2}
\begin{tabular}{ll}
\textbf{Exact datum} & \textbf{Derived question} \\
\hline
$\operatorname{gr}_F A\cong\mathbb C[J_\infty S_f]$ &
 convergence and collapse of the filtered chiral bar spectral sequence \\
generator-degree-$r$ OPE monomial &
 comparison with an arity-$r$ transferred operation \\
$c_{\mathrm{BP}}(k)$ from the KRW/FKR conformal vector &
 genus-$1$ modular curvature $\kappa_{\mathrm{BP}}(k)$ \\
transpose of the nilpotent partition &
 DS/bar transport and object-level Verdier companion
\end{tabular}
\end{center}

The Bershadsky--Polyakov algebra supplies the decisive worked case.  Its
four strong generators $J,G^+,G^-,T$ are even, its standard central charge is
\[
c_{\mathrm{BP}}(k)
=-\frac{(2k+3)(3k+1)}{k+3},
\qquad
c_{\mathrm{BP}}(k)+c_{\mathrm{BP}}(-k-6)=50,
\]
and its modular characteristic remains an open genus-$1$ computation.  The
secondary shifted conformal vector gives a different rational function whose
reflection sum is~$196$; that number belongs exclusively to the shifted lane
and remains separate from the standard FKR conductor.

\section{Completed bar--cobar from the PBW--Slodowy mechanism}

\begin{definition}[Filtered-complete chiral algebra]
Let $A$ be a chiral or vertex algebra equipped with an exhaustive separated
increasing filtration
\[
F_0A \subset F_1A \subset F_2A \subset \cdots \subset A.
\]
We assume that the filtration is multiplicative and compatible with the
translation operator $\partial$. The \emph{completed reduced bar complex}
$\widehat{\overline B}(A)$ is the completion of the usual reduced bar complex
with respect to the induced filtration by total $F$-degree.
\end{definition}

\begin{definition}[Completed Koszulity]
A filtered-complete chiral algebra $A$ is \emph{completed Koszul} if the
homology of $\widehat{\overline B}(A)$ is concentrated on the diagonal with
respect to the bar grading and the induced filtration grading, equivalently if
the completed cobar of this homology reconstructs $A$ up to filtered
quasi-isomorphism.
\end{definition}

\begin{theorem}[PBW--Slodowy input to the completed bar comparison;
\ClaimStatusConditional]
\label{thm:pbw-slodowy-collapse}
\textup{[}Type signature:
$(\mathrm{quadrant}=\mathrm{Open}\cap\mathrm{Chain},
\mathrm{presentation}=\mathrm{chain\text{-}level},
\mathrm{level}=1\leftrightarrow 2,
\mathrm{hypothesis\ package}=
\mathrm{exhaustive\ separated\ multiplicative\ filtration};
\mathrm{finite\text{-}dimensional\ weight\ slices};
\mathrm{Slodowy\text{-}slice\ free\ differential\ polynomial\ }\operatorname{gr}_F A;
H_{\mathrm{PBW}}^{\mathrm{bar}})$\textup{]}
Let $A$ be a filtered chiral algebra such that:
\begin{enumerate}[label=(\roman*)]
\item the filtration is exhaustive, separated, and multiplicative;
\item every conformal-weight piece is finite-dimensional and the filtration is
 bounded below on each weight;
\item there is an isomorphism of differential-commutative graded algebras
\[
\operatorname{gr}_F A \cong \operatorname{Sym}_{\partial}(V)
:= \operatorname{Sym}\!\left(\bigoplus_{m \geq 0}\partial^m V\right)
\]
for a finite-dimensional vector space $V$.
\end{enumerate}
Then the ordinary algebraic bar model of the associated graded has
\[
H_*\!\left(\overline B_{\mathrm{alg}}
(\operatorname{Sym}_{\partial}V)\right)
\cong \Lambda_{\partial}(sV).
\]
Here $H_{\mathrm{PBW}}^{\mathrm{bar}}$ denotes the additional package that
identifies the associated-graded chiral collision complex with this algebraic
bar model, proves complete convergence and strict Mittag--Leffler passage,
eliminates higher differentials and extension classes, and supplies the
completed twisting comparison.  Under
$H_{\mathrm{PBW}}^{\mathrm{bar}}$ the actual chiral bar spectral sequence has
$E_1^{\mathrm{PBW}}\cong\Lambda_{\partial}(sV)$ and one obtains
\[
H_*\!\left(\widehat{\overline B}(A)\right)
\cong \widehat{\Lambda}_{\partial}(sV),
\qquad
\widehat{\Omega}\,\widehat{\Lambda}_{\partial}(sV)
\xrightarrow{\sim} A.
\]
The first display is the algebraic PBW computation.  The chiral comparison
and completed Koszulity are the conclusions after the named package is
supplied.
\end{theorem}

\begin{proof}
Filter $\overline B(A)$ by total $F$-degree. Multiplicativity of the
filtration identifies the $E_0$-page with the reduced bar complex of the
associated graded:
\[
E_0 \cong \overline B(\operatorname{gr}_F A)
= \overline B(\operatorname{Sym}_{\partial}(V)).
\]
Therefore the $E_1$-page is the bar homology of the free differential polynomial
algebra on $V$. The latter is the exterior coalgebra on the suspended
derivative generators:
\[
H_*\!\left(\overline B(\operatorname{Sym}_{\partial}(V))\right)
\cong \Lambda_{\partial}(sV).
\]
$\operatorname{Sym}_{\partial}(V)$ is the tensor product of ordinary
polynomial algebras on the generators $\partial^m v$; the reduced bar homology
of an ordinary polynomial algebra is exterior on the suspension of the
generators, and the tensor product over all derivatives gives the stated result.

This proves the displayed algebraic bar calculation.  The first clause of
$H_{\mathrm{PBW}}^{\mathrm{bar}}$ transports it to the chiral collision
complex.  Its remaining clauses control higher differentials and the inverse
limit.  The complete filtered comparison theorem and the universal twisting
morphism then give the conditional bar-homology and cobar conclusions.
\end{proof}

\begin{corollary}[Principal affine \texorpdfstring{$W$}{W}-algebras:
PBW page and conditional completed Koszulity; \ClaimStatusConditional]
\label{cor:principal-w-completed-koszul}
\textup{[}Type signature:
$(\mathrm{Open}\cap\mathrm{Chain},\mathrm{chain\text{-}level},
\mathrm{level}=1\leftrightarrow 2,
\mathrm{hypothesis\ package}=\mathrm{simple\ }\fg;
f_{\mathrm{prin}}\mathrm{\ principal\ nilpotent};
k\ne -h^\vee;
\mathrm{Arakawa\ Kazhdan\text{-}filtration\ theorem\ for\ principal\ DS};
H_{\mathrm{PBW}}^{\mathrm{bar}})$\textup{]}
Let $\mathfrak g$ be simple and let $f_{\mathrm{prin}}$ be principal. Assume
$k\neq -h^\vee$. Then the universal affine $W$-algebra
$\mathcal W^k(\mathfrak g,f_{\mathrm{prin}})$ has free differential-polynomial
PBW associated graded.  Under $H_{\mathrm{PBW}}^{\mathrm{bar}}$, its chiral
bar spectral sequence has the exterior $E_1$ page and the algebra is completed
Koszul.
\end{corollary}

\begin{proof}
By Arakawa's theorem on the Li filtration in the principal case
\cite{Arakawa07}, the associated
graded of the universal affine $W$-algebra is the coordinate ring of the arc
space of the Kostant slice:
\[
\operatorname{gr}_{\mathrm{Li}}\mathcal W^k(\mathfrak g,f_{\mathrm{prin}})
\cong \mathbb C[J S_{\mathrm{prin}}].
\]
The Kostant slice is an affine space $f_{\mathrm{prin}}+\mathfrak g^e$, hence
its arc space is again affine and
\[
\mathbb C[J S_{\mathrm{prin}}]
\cong \operatorname{Sym}_{\partial}\bigl((\mathfrak g^e)^*\bigr).
\]
Apply Theorem~\ref{thm:pbw-slodowy-collapse} with
$V=(\mathfrak g^e)^*$.  The algebraic bar calculation follows directly; the
named package gives the chiral and completed conclusions.
\end{proof}

\begin{remark}[Conditional extension to arbitrary nilpotent data]
\label{rem:arbitrary-nilpotent-conditional}
For a general nilpotent~$f$, an associated-graded theorem of the form
\[
\operatorname{gr}_{\mathrm{Li}}\mathcal W^k(\mathfrak g,f)
\cong \mathbb C[J S_f],
\qquad
S_f=f+\mathfrak g^e.
\]
determines the PBW--Slodowy page.  Completed Koszulity additionally requires
$H_{\mathrm{PBW}}^{\mathrm{bar}}$ for the chosen chiral bar model.  The arc
space determines the associated graded; the quantum collision differential
and its completion determine the bar homology.
\end{remark}

\section{Reduction to hook-type families in type \texorpdfstring{$A$}{A}}

The completed theorem above solves the PBW side of the problem. The
representation-theoretic reduction of the \emph{classification} problem in type
$A$ is now supplied by inverse and staged Hamiltonian reduction. Taken
together, recent results of Fehily, Creutzig--Linshaw--Nakatsuka--Sato,
Creutzig--Fasquel--Linshaw--Nakatsuka, Genra--Juillard, and Butson--Nair show
that generic type-$A$ affine $W$-algebras can be related along the
nilpotent-orbit order by successive reduction and inverse reduction steps, and
that the hook-type locus forms the natural irreducible spine of this network
\cite{FehilyHook,CLNS24,CFLN24,GenraStages,ButsonNair}. Concretely, inverse
reductions exist between the hook-type $\mathfrak{sl}_{n+1}$ $W$-algebras;
reduction by stages is proved under compatibility conditions; and inverse
Hamiltonian reduction for all affine $W$-algebras in type~$A$ at generic
level is now available. These results supply the reduction-theoretic spine.
The Koszul-duality statement along that spine is the conditional
DS--bar-compatibility corridor of
Theorem~\ref{thm:hook-transport-corridor}.

\begin{principle}[Transport/orbit factorization]
\label{princ:transport-orbit-factorization}
\index{W-algebra@$\mathcal{W}$-algebra!transport/orbit factorization}
The non-principal duality problem splits into two axes:
\[
\text{non-principal duality}
\;=\;
\text{transport along reduction network}
\;+\;
\text{identification of dual orbit/level}.
\]
The reduction-network axis is proved as a transport graph on the
hook spine and through the verified closure ranges; the
Koszul-duality transport along that axis remains conditional on
DS--bar compatibility. The second axis remains the true frontier.
\end{principle}

\begin{definition}[Reduction graph]
\label{def:reduction-graph}
\ClaimStatusDefinitional
\index{reduction graph!$\Gamma_N$}
Let $\Gamma_N$ be the graph whose vertices are the partitions
$\lambda \vdash N$ (indexing nilpotent orbits in
$\mathfrak{sl}_N$), and whose edges are the proved
reduction or inverse-reduction functors between the corresponding
type-$A$ affine $\mathcal{W}$-algebras at generic level.
The \emph{hook vertices} are the partitions of hook type
$(N-r, 1^r)$, $0 \leq r \leq N-1$.
The \emph{transport-closure} of a set $S$ of vertices is the set
of all vertices reachable from $S$ by paths in $\Gamma_N$.
\end{definition}

\begin{theorem}[Hook-type transport corridor under DS--bar compatibility; \ClaimStatusConditional]
\label{thm:hook-transport-corridor}
\index{W-algebra@$\mathcal{W}$-algebra!hook transport corridor}
\textup{[}Type signature:
$(\mathrm{Open}\cap\mathrm{Cat},\mathrm{factorisation},
\mathrm{level}=1\leftrightarrow 2,
\mathrm{hypothesis\ package}=\mathrm{generic\ level\ }k;
\mathrm{Fehily\ inverse\text{-}Hamiltonian\ reduction\ for\ hooks};
\mathrm{Creutzig\text{--}Linshaw\text{--}Nakatsuka\text{--}Sato\ \textup{(}CLNS\textup{)} hook\ duality\ at\ generic\ }k;
\mathrm{Genra\ reduction\ by\ stages\ at\ admissible\ }k;
\mathrm{Butson\text{--}Nair\ geometric\ inverse\ Hamiltonian\ reduction};
H_{\mathrm{hook}}^{\mathrm{DS/bar}};
\mathrm{finite\text{-}type\ Verdier\ dual\ or\ completed\ continuous\ nondegenerate\ pairing})$\textup{]}
Assume that, at generic level, bar--cobar/Koszul duality intertwines
with the reduction and inverse-reduction functors along the hook
network in type~$A$. Then for every hook partition
$\eta = (N-r, 1^r)$ in $\mathfrak{sl}_N$:
\[
\mathcal{W}^k(\mathfrak{sl}_N, f_\eta)^!
\;\simeq\;
\mathcal{W}^{-k-2N}(\mathfrak{sl}_N, f_{\eta^t}).
\]
For a self-transpose hook, the formula selects a same-family candidate at the
reflected level.  Membership in KSDual is the stronger assertion that the
Verdier comparison is an object-level fixed-point equivalence with a
nondegenerate involutive pairing.  It is an additional clause, rather than a
consequence of $\eta=\eta^t$.
\end{theorem}

\begin{proof}
The cited reduction and inverse-reduction results construct the hook network.
The package $H_{\mathrm{hook}}^{\mathrm{DS/bar}}$ asserts that each edge lifts
to the completed bar object, commutes with Verdier duality, and identifies the
transposed target.  Induction along a path then transports the principal
Feigin--Frenkel comparison to the stated candidate.  This proof is conditional
exactly at the edge-lifting step.
\end{proof}

\begin{proposition}[Transport propagation lemma; \ClaimStatusProvedHere]
\label{prop:transport-propagation}
\index{transport propagation}
Suppose a candidate duality functor $D$ satisfies:
\begin{enumerate}[label=\textup{(\roman*)}]
\item $D\bigl(\mathcal{W}^k(\mathfrak{sl}_N, f_\eta)\bigr)
 \simeq \mathcal{W}^{k_\eta^\vee}(\mathfrak{sl}_N, f_{\eta^t})$
 for every hook partition~$\eta$, and
\item $D$ intertwines every edge of $\Gamma_N$ with the
 transposed edge.
\end{enumerate}
Then $D$ is forced to satisfy
$D\bigl(\mathcal{W}^k(\mathfrak{sl}_N, f_\lambda)\bigr)
\simeq \mathcal{W}^{k_\lambda^\vee}(\mathfrak{sl}_N, f_{\lambda^t})$
for every partition $\lambda$ in the transport-closure of the
hook vertices.
\end{proposition}

\begin{proof}
Pick a path in $\Gamma_N$ from a hook vertex to~$\lambda$.
Start from the known hook formula and propagate inductively across
the path using edge compatibility~(ii).
\end{proof}

\begin{remark}[Formality mechanism in the hook corridor]
\label{rem:hook-formality-mechanism}
\index{hook transport corridor!formality mechanism}
The hook-type corridor (Theorem~\ref{thm:hook-transport-corridor})
is conjectured to follow from the formality route
(Proposition~\ref{prop:ds-bar-formality}). This requires uniform
DS/bar compatibility across the hook network.  Abelian
$\mathfrak n_+$ cases support the filtration-formality mechanism case by
case.  Hook-wide transport remains a conditional corridor, and $(3,2)$ in
$\mathfrak{sl}_5$ is the
first explicit non-abelian audit surface
(Remark~\ref{rem:abelian-nonabelian-nilradical}).
\end{remark}

\begin{conjecture}[Type-$A$ transport-to-transpose; \ClaimStatusConjectured]
\label{conj:type-a-transport-to-transpose}
\index{transport-to-transpose conjecture}
At generic level, the transport-closure of the hook vertices is
all of $\operatorname{Par}(N)$, and Koszul/bar--cobar duality
intertwines every proved reduction or inverse-reduction edge.
Consequently,
\[
\bigl(\mathcal{W}^k(\mathfrak{sl}_N, f_\lambda)\bigr)^!
\;\simeq\;
\mathcal{W}^{k_\lambda^\vee}(\mathfrak{sl}_N, f_{\lambda^t}).
\]
\end{conjecture}

\begin{remark}[Relation to the W-orbit duality conjecture]
Conjecture~\ref{conj:type-a-transport-to-transpose} is the
type-$A$ specialization of Conjecture~\ref{conj:w-orbit-duality},
with the global miracle replaced by an edge-compatibility
theorem. It is sharper than the old ``arbitrary BV duality''
slogan (it identifies the precise mechanism, transport through
$\Gamma_N$) but still large enough to deserve the frontier label.
\end{remark}

\subsection{Hook-type ghost constants and the DS--bar commutation}
\label{sec:hook-ghost-ds-bar}
\index{W-algebra@$\mathcal{W}$-algebra!hook ghost constant}
\index{DS reduction!bar commutation}

The ghost constant $C_\lambda$ measures the conformal-weight
contribution of the ghost system in DS reduction.
For hook partitions this admits a closed formula.

\begin{proposition}[Hook ghost constant; \ClaimStatusProvedHere]
\label{prop:hook-ghost-constant}
\index{ghost constant!hook type}
For $\fg = \mathfrak{sl}_N$ with hook partition
$\lambda = (m, 1^r)$, $m + r = N$, the ghost constant is
\begin{equation}\label{eq:hook-ghost-constant}
C_{(m,1^r)} \;=\; \frac{m(m^2-1)}{6}
 + \frac{r \lfloor m^2/2 \rfloor}{2}.
\end{equation}
\end{proposition}

\begin{proof}
The ghost system $\mathcal{F}_{\mathrm{gh}}(\mathfrak{m}^+)$
is indexed by the positive eigenspaces of $\operatorname{ad}(h)/2$
on the nilradical $\mathfrak{m}^+$. For a hook partition
$(m,1^r)$, the $\mathfrak{sl}_2$-triple decomposes $\fg^e$ into
representations whose conformal-weight contributions to the
ghost central charge are computed from the partition combinatorics.
The first term $m(m^2-1)/6$ is the principal contribution
(ghost constant of $\mathcal{W}_m$); the second term accounts for
the $r$ additional generators introduced by the hook tail.
\end{proof}

\begin{computation}[$\mathfrak{sl}_4$ hook-type data]
\label{comp:sl4-hook-data}
\index{W-algebra@$\mathcal{W}$-algebra!sl4 data@$\mathfrak{sl}_4$ data}
\begin{center}
\begin{tabular}{llcl}
\toprule
Partition $\lambda$ & $\dim(\fg^{f_\lambda})$ & $C_\lambda$
 & Strong generators \\
\midrule
$(4)$ principal & 3 & 10
 & $W_2, W_3, W_4$ ($= \mathcal{W}_4$ algebra) \\
$(3,1)$ subregular & 5 & 6
 & $1 \times h\!=\!1,\; 3 \times h\!=\!2,\; 1 \times h\!=\!3$ \\
$(2,2)$ even & 7 & 4
 & $3 \times h\!=\!1,\; 4 \times h\!=\!2$ \\
$(2,1,1)$ minimal & 9 & 3
 & $4 \times h\!=\!1,\; 4 \times h\!=\!\tfrac32,\; 1 \times h\!=\!2$ \\
\bottomrule
\end{tabular}
\end{center}
The hook partitions $(4)$, $(3,1)$, $(2,1,1)$ satisfy
\eqref{eq:hook-ghost-constant}. The even partition $(2,2)$ lies outside the
hook family and is included as the first rectangular audit surface.
\end{computation}

The DS presentation supplies the following data before any modular-curvature
comparison is attempted.

\begin{proposition}[Hook-type data visible before DS/bar comparison;
\ClaimStatusProvedElsewhere]
\label{prop:ds-bar-hook-commutation}
\index{DS reduction!bar commutation!hook type}
For every hook partition $\eta = (N-r, 1^r)$ in
$\mathfrak{sl}_N$ at generic level~$k$, the DS reduction
$\mathcal{W}^k(\mathfrak{sl}_N, f_\eta)$ carries the following
reduction-theoretic data by Kac--Roan--Wakimoto~\cite{KRW03} and
Arakawa~\cite{Arakawa07}:
\begin{enumerate}[label=\textup{(\roman*)}]
\item \emph{Generator matching}: the strong generating set of
 $\mathcal{W}^k(\mathfrak{sl}_N, f_\eta)$ has
 $\dim(\mathfrak{sl}_N^{f_\eta})$ elements, coinciding with the
 BRST-surviving generators of the DS complex.
\item \emph{Central charge}: the Kac--Roan--Wakimoto formula computes
 the central charge for the chosen good grading and conformal vector.
\item \emph{Filtered classical limit}: the Kazhdan filtration identifies the
 associated graded BRST cohomology with the arc algebra of the Slodowy slice
 on the regular universal lane covered by the cited DS theorem.
\end{enumerate}
The modular characteristic, Theorem~H, completed bar collapse, and
commutation of DS with Verdier/bar duality require independent comparison
theorems.  In particular, the reciprocal-weight signed sum of a strong
generating set is an OPE-spectrum diagnostic.  The genus-$1$ curvature gives
the independent definition of~$\kappa$.
\end{proposition}

\begin{proof}
The generator count and filtered classical limit are the standard BRST and
Kazhdan-filtration conclusions for affine $W$-algebras.  The central charge is
the KRW supertrace formula evaluated in the same conformal-vector convention.
These statements precede any construction of a chiral bar comparison map.
\end{proof}

\begin{remark}[Transport-to-transpose scalar audit]
\label{rem:hook-kappa-sum-rule}
\index{kappa@$\kappa$!transport sum rule}
The KRW central-charge functions can be tested under the candidate reflection
$k\mapsto-k-2N$ and orbit transpose.  A level-independent sum is exact
evidence on the conformal-anomaly lane.  A modular complementarity identity
requires the genus-$1$ curvature on both sides and the DS/bar/Verdier package.
For BP the exact standard central-charge sum is~$50$; its modular
complementarity sum remains open.
\end{remark}

\section{OPE generator degree and the derived comparison problem}

\begin{definition}[Generator degree]
Fix an ordered strong generating set $V=\langle X_1,\dots,X_N\rangle$ for a
filtered-complete chiral algebra $A$. The \emph{generator degree} of a
normally ordered differential monomial is the number of generators $X_i$ that
appear in it; derivatives carry degree~$0$.  The maximal generator
degree among the singular OPE coefficients is denoted
$d_{\mathrm{OPE}}(A;V)$.  It is an invariant of the chosen normal form.
\end{definition}

\begin{theorem}[Generator-degree detection in an OPE normal form;
\ClaimStatusProvedHere]
\label{thm:canonical-degree-detection}
Let $A$ be presented in a fixed normal form by strong generators
$X_1,\dots,X_N$. Suppose that every singular OPE coefficient among these
generators is a normally ordered differential polynomial of generator degree at
most $d$. Then $d_{\mathrm{OPE}}(A;V)\le d$.  If one coefficient contains a
nonzero degree-$d$ term in that normal form, then
$d_{\mathrm{OPE}}(A;V)=d$.
\end{theorem}

\begin{proof}
This is the filtration of the normally ordered differential-polynomial algebra
by the number of strong generators.  The upper bound and the equality follow
directly from the definition.
\end{proof}

\begin{remark}[The comparison required for a bar-dual statement]
An identification between generator degree and a Taylor coefficient of a
minimal or completed bar dual requires an explicit contraction, a proof that
the chiral collision differential is represented by the chosen OPE
polynomials, and a non-cancellation statement under homotopy transfer.  These
data form a comparison package $H_{\mathrm{OPE/bar}}$.  The calculations below
establish exact OPE degrees.  Their promotion to $A_\infty$ arities is an open
problem under $H_{\mathrm{OPE/bar}}$.
\end{remark}

\section{Screening-fiber duality for inverse reduction}

\begin{construction}[Screening-fiber dual]
\label{constr:screening-fiber-dual}
Suppose that a hook-type affine $W$-algebra $A$ is realized as a joint kernel of
screening operators inside a larger algebra:
\[
A \simeq \ker(S_1,\dots,S_r)\subset B\otimes \mathcal F,
\]
where $B$ is another affine $W$-algebra, $\mathcal F$ is a free-field algebra,
and the standard screening complex resolves $A$. Define the
\emph{screening-fiber dual} by
\[
\mathfrak D(A):=
\Bigl(
B^!\,\widehat\otimes\,\mathcal F^!\,\widehat\otimes\,
\Lambda(\eta_1,\dots,\eta_r),
\;
d_{B^!}+d_{\mathcal F^!}+\sum_{i=1}^r \eta_i\,S_i^\vee
\Bigr).
\]
\end{construction}

\begin{conjecture}[Screening-fiber model for the completed dual; \ClaimStatusConjectured; \textnormal{[conditional on screening exactness]}]
\label{conj:screening-fiber-dual}
Under the usual exactness and convergence hypotheses for the screening
resolution, $\mathfrak D(A)$ is a strict dg model whose transferred
$A_\infty$-minimal model computes the completed dual of $A$.
\end{conjecture}

\begin{remark}[Screening transfer criterion]
The conjecture is equivalent to the following chain-level packet.
First, the screening resolution
\[
A\longrightarrow (B\otimes\mathcal F,d_{\mathrm{scr}})
\]
must be exact in the completed category used for the reduced bar
construction. Second, the completed bar construction must commute with
the screening fibre, so that the bar-dual coalgebra of~$A$ is the
homology of the completed fibre of
\[
\barB(B\otimes\mathcal F)
\xrightarrow{\;\barB(S_1,\ldots,S_r)\;}
\bigoplus_i\barB(B\otimes\mathcal F).
\]
Third, continuous Verdier duality must send this fibre to the displayed
dg algebra
\[
B^!\widehat\otimes \mathcal F^!
\widehat\otimes \Lambda(\eta_1,\ldots,\eta_r),
\qquad
d=d_{B^!}+d_{\mathcal F^!}+\sum_i\eta_iS_i^\vee.
\]
Under these three hypotheses, homological perturbation transfers this
strict dg algebra to a minimal $A_\infty$ model for the completed
dual. The higher multiplications are then precisely the Massey
products created by the screening differential; the free-field factor remains
on its strict quadratic lane.
\end{remark}

This construction is the homological content of inverse reduction: higher
operations in hook-type duals are screening Massey products.

\section{The Feigin--Semikhatov realization and the exact coefficient packet}

The maximally asymmetric realization of the subregular family
$\mathcal W_n^{(2)}$ is due to Feigin--Semikhatov
\cite{FeiginSemikhatov}, write
\[
E_n(z)=e^{\Xi}(z),
\qquad
F_n(z)=-P_n^{(k)}(A_{n-1},\dots,A_1,Q)\,e^{-\Xi}(z),
\]
where
\[
P_n^{(k)}
=
\bigl((k+n-1)\partial + Q+A_{n-1}+\cdots+A_1\bigr)\circ\cdots\circ
\bigl((k+n-1)\partial + Q+A_1\bigr)\circ Q.
\]
Set
\[
K_n:=k+n-1,
\qquad
u_j:=Q+A_1+\cdots+A_j
\quad (0\leq j\leq n-1),
\]
with the convention $u_0=Q$.

\subsection{The exact raw singular coefficients}

The crucial point is that one can write \emph{every} singular coefficient in
the $E_n$--$F_n$ OPE recursively.

\begin{definition}[Recursive coefficient operators; \ClaimStatusDefinitional]
\label{def:raw-coefficient-operators}
For fixed $n$, define differential polynomials
$\mathcal C_{j,m}=\mathcal C_{j,m}^{(n,k)}$ recursively by
\[
\mathcal C_{0,0}=1,
\qquad
\mathcal C_{j,m}=0 \;\text{ if }\; m<0 \;\text{ or }\; m>j,
\]
and
\[
\mathcal C_{j+1,m}
=
\bigl(K_n(j-m)-1\bigr)\,\mathcal C_{j,m}
+
\bigl(K_n\partial + u_j\bigr)\mathcal C_{j,m-1}
\qquad
(0\leq m\leq j+1).
\]
\end{definition}

\begin{definition}[Bell polynomials of the exponential factor; \ClaimStatusDefinitional]
\label{def:bell-polynomials-exp}
Define differential polynomials $B_a(Y,\partial Y,\dots,\partial^{a-1}Y)$ by
\[
\exp\!\left(
\sum_{r\geq 0}\frac{t^{r+1}}{(r+1)!}\partial^r Y
\right)
=
\sum_{a\geq 0}\frac{t^a}{a!}\,B_a(Y,\partial Y,\dots,\partial^{a-1}Y),
\]
so that
\[
e^{\Xi(z)-\Xi(w)}
=
\sum_{a\geq 0}
\frac{(z-w)^a}{a!}\,
B_a(Y,\partial Y,\dots,\partial^{a-1}Y)(w).
\]
\end{definition}

\begin{theorem}[Exact Bell recursion for the full singular packet; \ClaimStatusProvedHere]
\label{thm:full-raw-coefficient-packet}
Let $t=z-w$. In the maximally asymmetric realization of
$\mathcal W_n^{(2)}(k)$ one has
\[
E_n(z)F_n(w)
=
\sum_{r=0}^{n-1}\frac{U_{n,r}^{\mathrm{raw}}(w)}{t^{n-r}}
+ O(t^0),
\]
where the singular coefficients are
\[
U_{n,r}^{\mathrm{raw}}
=
-
\sum_{a=0}^{r}
\frac{1}{a!}\,
B_a(Y,\partial Y,\dots,\partial^{a-1}Y)\,
\mathcal C_{n,r-a}.
\]
Equivalently, the entire singular part of the $E_n$--$F_n$ OPE is determined
by the Bell polynomials of the exponential factor and the recursion
\[
\mathcal C_{j+1,m}
=
\bigl(K_n(j-m)-1\bigr)\,\mathcal C_{j,m}
+
\bigl(K_n\partial + u_j\bigr)\mathcal C_{j,m-1}.
\]
\end{theorem}

\begin{proof}
Feigin--Semikhatov write the maximally asymmetric generator as
\[
F_n(w)=-\bigl((K_n\partial + u_{n-1})\circ\cdots\circ(K_n\partial+u_0)\bigr)e^{-\Xi}(w).
\]
Move $e^{\Xi}(z)$ past the ordered product by Wick expansion. Since
$\partial_w e^{-\Xi}(w)$ produces $Y(w)e^{-\Xi}(w)$ and
$Q(w)e^{\Xi}(z)$ contributes $-(z-w)^{-1}e^{\Xi}(z)$, each factor
$K_n\partial_w+u_j(w)$ is replaced by the ordered operator
\[
K_n\partial_w + u_j(w) - t^{-1}
\]
acting on everything to its right. Thus
\[
E_n(z)F_n(w)
=
-\,e^{\Xi(z)-\Xi(w)}\,
\Bigl[(K_n\partial_w+u_{n-1}-t^{-1})\cdots(K_n\partial_w+u_0-t^{-1})\cdot 1\Bigr].
\]
Define
\[
R_j(t):=(K_n\partial_w+u_{j-1}-t^{-1})\cdots(K_n\partial_w+u_0-t^{-1})\cdot 1
\qquad (j\geq 1),
\]
with $R_0(t)=1$. Expand
\[
R_j(t)=\sum_{m=0}^j \mathcal C_{j,m}\,t^{-(j-m)}.
\]
Applying one more factor $K_n\partial_w+u_j-t^{-1}$ and using
\[
\partial_w \bigl(t^{-p}\bigr)= p\,t^{-(p+1)}
\qquad (t=z-w)
\]
gives exactly the recursion of
Definition~\ref{def:raw-coefficient-operators}. This proves
\[
R_n(t)=\sum_{m=0}^n \mathcal C_{n,m}\,t^{-(n-m)}.
\]

Now expand the exponential factor as in
Definition~\ref{def:bell-polynomials-exp}:
\[
e^{\Xi(z)-\Xi(w)}
=
\sum_{a\geq 0}
\frac{t^a}{a!}\,
B_a(Y,\partial Y,\dots,\partial^{a-1}Y).
\]
Multiply by the expansion of $R_n(t)$ and collect the coefficient of
$t^{-(n-r)}$. The result is the displayed formula for
$U_{n,r}^{\mathrm{raw}}$.
\end{proof}

\begin{remark}[Why this answers the coefficient problem]
Theorem~\ref{thm:full-raw-coefficient-packet} writes \emph{every} singular
coefficient in the $E_n$--$F_n$ OPE exactly, including all lower-derivative
terms. Since the other basic OPEs are fixed by the Heisenberg and Virasoro
requirements,
\[
H_n(z)H_n(w)\sim \frac{\ell_n(k)}{(z-w)^2},
\qquad
H_n(z)E_n(w)\sim \frac{E_n(w)}{z-w},
\qquad
H_n(z)F_n(w)\sim -\frac{F_n(w)}{z-w},
\]
and the higher neutral generators are extracted as the pole coefficients of
$E_n(z)F_n(w)$, the full VOA is reconstructed from these raw coefficient
packets together with primarity and the Jacobi identities.
\end{remark}

\subsection{Low-order checks}

The recursion immediately recovers the known leading coefficients:
$\mathcal C_{n,0}=(-1)^n\lambda_{n-1}(n,k)$, where
\[
\lambda_m(n,k)=\prod_{i=1}^{m}\bigl(i(k+n-1)-1\bigr),
\]
so the leading pole is
\[
U_{n,0}^{\mathrm{raw}}=\lambda_{n-1}(n,k).
\]
Likewise one recovers the dimension-$1$ pole
\[
U_{n,1}^{\mathrm{raw}}=n\lambda_{n-2}(n,k)H_n,
\]
and, after reorganizing the weight-$2$ coefficient into Heisenberg and Virasoro
pieces, the Feigin--Semikhatov formula
\cite[(A.1)]{FeiginSemikhatov}:
\begin{align}
E_n(z)F_n(w)
&=
\frac{\lambda_{n-1}(n,k)}{(z-w)^n}
+
\frac{n\lambda_{n-2}(n,k)H_n(w)}{(z-w)^{n-1}}
\notag\\[4pt]
&\quad
+\frac{\lambda_{n-3}(n,k)\left(\frac{n(n-1)}2 H_nH_n(w)
+\frac{n((n-2)(k+n-1)-1)}2\,\partial H_n(w)
-(k+n)T_n(w)\right)}{(z-w)^{n-2}}
+\cdots.
\label{eq:fs-pole-two}
\end{align}

\section{The subregular Miura--Appell symbols}

The exact recursion above controls every singular coefficient. To isolate the
highest generator-degree part one passes to a derivative-killing symbol map.

\begin{definition}[Derivative-killing symbol map]
\label{def:derivative-killing-symbol}
Let
\[
\sigma:\ \mathbb C[Q,A_1,\dots,A_{n-1},Y,\partial Q,\partial A_i,\partial Y,\dots]
\longrightarrow
\mathbb C[q,a_1,\dots,a_{n-1},y]
\]
be the algebra map killing all derivatives and sending
$Q\mapsto q$, $A_i\mapsto a_i$, $Y\mapsto y$.
\end{definition}

In the maximally asymmetric realization only the combination
\[
u_j:=q+a_1+\cdots+a_j
\qquad (0\leq j\leq n-1)
\]
enters the symbol of $P_n$.

\begin{theorem}[Subregular Miura product formula; \ClaimStatusProvedHere]
\label{thm:miura-product-formula}
With $u_j$ as above,
\[
\sigma\!\left(P_n^{(k)}\right)=\prod_{j=0}^{n-1}u_j.
\]
\end{theorem}

\begin{proof}
Feigin--Semikhatov's recursion reads
\[
P_n^{(k)}=\bigl((k+n-1)\partial + u_{n-1}\bigr)P_{n-1}^{(k+1)},
\qquad
P_1^{(k)}=Q.
\]
Applying $\sigma$ kills the derivative term and yields
\[
\sigma(P_n^{(k)})=u_{n-1}\,\sigma(P_{n-1}^{(k+1)}).
\]
Since $\sigma(P_1^{(k)})=u_0=q$, induction proves the claim.
\end{proof}

Introduce the center-of-mass coordinate
\[
h_n:=\frac{1}{n}\sum_{j=0}^{n-1}u_j
=
q+\frac{1}{n}\sum_{i=1}^{n-1}(n-i)a_i
\]
and the centered variables
\[
v_j:=u_j-h_n,
\qquad
\sum_{j=0}^{n-1}v_j=0.
\]
Let $e_r(v)$ denote the elementary symmetric polynomial of degree $r$ in
$v_0,\dots,v_{n-1}$.

\begin{theorem}[Subregular Appell formula; \ClaimStatusProvedHere]
\label{thm:subregular-appell-formula}
For $2\leq r\leq n-1$, the leading symbol of the weight-$r$ raw coefficient is
\[
C_{n,r}(h_n,v)
:=
\frac{1}{(n-r)!}\,\partial_q^{\,n-r}\!\left(\prod_{j=0}^{n-1}u_j\right)
=
\sum_{t=0}^{r}\binom{n-r+t}{t}\,e_{r-t}(v)\,h_n^t.
\]
Equivalently,
\[
\partial_{h_n}C_{n,r}=(n-r+1)\,C_{n,r-1}.
\]
\end{theorem}

\begin{proof}
Since $\prod_{j=0}^{n-1}u_j=\prod_{j=0}^{n-1}(h_n+v_j)$ and the $v_j$ are
independent of $q$ while $\partial_q h_n=1$, one has
\[
\frac{1}{(n-r)!}\partial_q^{\,n-r}
=
\frac{1}{(n-r)!}\partial_{h_n}^{\,n-r}.
\]
Expand the product in elementary symmetric functions:
\[
\prod_{j=0}^{n-1}(h_n+v_j)=\sum_{s=0}^{n}e_{n-s}(v)\,h_n^{\,s}.
\]
Differentiate termwise. Writing $s=n-r+t$ gives
\[
\frac{1}{(n-r)!}\partial_{h_n}^{\,n-r}h_n^{\,s}
=
\binom{s}{n-r}h_n^{\,s-n+r}
=
\binom{n-r+t}{t}h_n^t.
\]
Substituting into the expansion proves the formula.
\end{proof}

\begin{remark}[Normalized symbols]
The centered variable $h_n$ carries the Heisenberg direction, while the
translation-invariant part $e_r(v)$ is the normalized symbol of the new neutral
generator in weight $r$. Thus the Appell formula separates the Heisenberg
tower from the orthogonal, $H$-regular generator at the symbolic level.
\end{remark}

\section{The first non-principal stages}

\subsection{The Bershadsky--Polyakov algebra}

\begin{proposition}[Bershadsky--Polyakov OPEs in Feigin--Semikhatov
normal form; \ClaimStatusProvedElsewhere]
\label{prop:bp-fs-normal-form-opes}
For $n=3$ one recovers the Bershadsky--Polyakov algebra. Feigin--Semikhatov
write its OPEs as follows \cite[(A.3)]{FeiginSemikhatov}:
\begin{align}
E(z)F(w)
&=
\frac{(k+1)(2k+3)}{(z-w)^3}
+
\frac{3(k+1)H(w)}{(z-w)^2}
+
\frac{3HH(w)-(k+3)T(w)+\frac32(k+1)\partial H(w)}{z-w},
\label{eq:bp-ef}
\\[4pt]
H(z)H(w)
&\sim \frac{2k+3}{3(z-w)^2},
\qquad
H(z)E(w)\sim \frac{E(w)}{z-w},
\qquad
H(z)F(w)\sim -\frac{F(w)}{z-w}.
\label{eq:bp-heis}
\end{align}
All remaining OPEs are linear or quadratic in the generators.
\end{proposition}

\begin{theorem}[Quadratic OPE degree of Bershadsky--Polyakov;
\ClaimStatusProvedHere]
\label{thm:bp-strict}
\textup{[}Type signature:
$(\mathrm{Open}\cap\mathrm{Chain},\mathrm{chain\text{-}level},
\mathrm{level}=1\leftrightarrow 2,
\mathrm{hypothesis\ package}=k\ne -3;
\mathrm{Feigin\text{--}Semikhatov\ realisation};
\mathrm{OPE\ normal\ form\ at\ generators\ }J,G^\pm,T)$\textup{]}
For every $k\neq-3$, the Feigin--Semikhatov generator presentation of
$\mathcal W_3^{(2)}(k)$ has
\[
d_{\mathrm{OPE}}(\mathcal W_3^{(2)};J,G^+,G^-,T)=2.
\]
The quadratic term is $3{:}JJ{:}$ in the simple pole of
$G^+(z)G^-(w)$.
\end{theorem}

\begin{proof}
The singular OPE coefficients in
\eqref{eq:bp-ef}--\eqref{eq:bp-heis} have generator degree at most~$2$, and
the displayed $HH$ coefficient is nonzero.  Apply
Theorem~\ref{thm:canonical-degree-detection}.
\end{proof}

\begin{computation}[Bershadsky--Polyakov scalar audit;
\ClaimStatusProvedHere]
\label{comp:bp-kappa-three-paths}
\index{Bershadsky--Polyakov!modular characteristic}
\index{kappa@$\kappa$!Bershadsky--Polyakov}
Fehily--Kawasetsu--Ridout~\cite[Definition~2.1 and
Eq.~\textup{(}2.2\textup{)}]{FKR20} use the ordinary bosonic vertex algebra
with even generators $J,G^+,G^-,T$ and standard central charge
\begin{equation}
\label{eq:bp-standard-central-audit}
c_{\mathrm{BP}}(k)
=-\frac{(2k+3)(3k+1)}{k+3}.
\end{equation}
Writing $t=k+3$ gives
\[
c_{\mathrm{BP}}(k)=25-6t-\frac{24}{t},
\qquad
c_{\mathrm{BP}}(-k-6)=25+6t+\frac{24}{t},
\]
and hence
\begin{equation}
\label{eq:bp-standard-conductor-audit}
c_{\mathrm{BP}}(k)+c_{\mathrm{BP}}(-k-6)=50.
\end{equation}
The midpoint value~$25$ occurs at the regular fibres $k=-3\pm2i$;
$k=-3$ is the common pole.

The secondary shifted conformal vector has
\begin{equation}
\label{eq:bp-shifted-secondary-audit}
c_{\mathrm{BP}}^{\mathrm{shift}}(k)
=2-\frac{24(k+1)^2}{k+3},
\qquad
c_{\mathrm{BP}}^{\mathrm{shift}}(k)
+c_{\mathrm{BP}}^{\mathrm{shift}}(-k-6)=196.
\end{equation}
Equations~\eqref{eq:bp-standard-central-audit} and
\eqref{eq:bp-shifted-secondary-audit} belong to different conformal-vector
conventions.

The unsigned reciprocal-weight sum of the even strong generators is
\[
1+\frac23+\frac23+\frac12=\frac{17}{6}.
\]
It records only the strong-generator spectrum.  The modular characteristic
$\kappa_{\mathrm{BP}}(k)$ is the coefficient of the non-separating genus-$1$
curvature and is \ClaimStatusOpen{} pending a direct calculation including
charged ghosts, neutral fields, the improvement term, and mixed channels.
\end{computation}

\begin{proposition}[Standard Bershadsky--Polyakov central-charge conductor;
\ClaimStatusProvedHere]
\label{prop:bp-complementarity-constant}
\index{Bershadsky--Polyakov!complementarity constant}
\index{central-charge conductor!Bershadsky--Polyakov}
In the standard FKR conformal-vector convention,
\[
K^c_{\mathrm{BP}}
:=c_{\mathrm{BP}}(k)+c_{\mathrm{BP}}(-k-6)=50.
\]
The shifted secondary convention of
\eqref{eq:bp-shifted-secondary-audit} has reflection sum~$196$.
\end{proposition}

\begin{proof}
This is Equation~\eqref{eq:bp-standard-conductor-audit}.  The second
statement is Equation~\eqref{eq:bp-shifted-secondary-audit}.
\end{proof}

\begin{proposition}[Bershadsky--Polyakov subregular ledger; \ClaimStatusConditional]
\label{prop:bp-subregular-standard-family-ledger}
\index{Bershadsky--Polyakov!standard family ledger}
\index{Drinfeld--Sokolov reduction!subregular}
\textup{[}Type signature:
$(\mathrm{Open}\cap\mathrm{Chain},\mathrm{chain\text{-}level},
\mathrm{level}=1,
\mathrm{hypothesis\ package}=k\ne -3;
\mathrm{FKR\ standard\ conformal\ vector};
\mathrm{Feigin\text{--}Semikhatov\ OPE\ normal\ form})$\textup{]}
The Bershadsky--Polyakov algebra is the $n=3$ member of the subregular
Drinfeld--Sokolov family. In the Feigin--Semikhatov normal form its strong
generators are
\begin{equation}
\label{eq:bp-ledger-fields}
H_{h=1}^{\mathrm{even}},\qquad
E_{h=3/2}^{\mathrm{even}},\qquad
F_{h=3/2}^{\mathrm{even}},\qquad
T_{h=2}^{\mathrm{even}},
\end{equation}
with OPEs \eqref{eq:bp-ef}--\eqref{eq:bp-heis}. Its central charge is
\begin{equation}
\label{eq:bp-ledger-central-charge}
c_{\mathrm{BP}}(k)
=-\frac{(2k+3)(3k+1)}{k+3},
\qquad k\ne -3,
\end{equation}
with exact reflection identity
\begin{equation}
\label{eq:bp-ledger-conductors}
K^c(\mathrm{BP}_k,\mathrm{BP}_{-k-6})=50.
\end{equation}
The shifted scalar $c_{\mathrm{BP}}^{\mathrm{shift}}$ has conductor~$196$
and is recorded separately in
\eqref{eq:bp-shifted-secondary-audit}.  The quantities
$\kappa_{\mathrm{BP}}$, $\varrho_{\mathrm{BP}}$, and
$K^\kappa_{\mathrm{BP}}$ are open genus-$1$ invariants.

The BP instance of Theorem~H has \ClaimStatusConditional{} under the named
package
\[
H_{\mathrm{BP}}^{H}:=
H_{\mathrm{PBW}}^{\mathrm{bar}}
+\{\text{finite/perfect Verdier duality, strict ML completion,
ordered acyclicity, averaging descent}\}.
\]
This package is the required bridge from finite OPE calculations to
curve-level chiral Hochschild amplitude and growth.

The shadow obstruction lines must be read separately:
\begin{equation}
\label{eq:bp-ledger-shadow-lines}
T\text{-line}:\mathbf M,\ r_{\max}=\infty;\qquad
H\text{-line}:\mathbf G,\ r_{\max}=2;\qquad
E,F\text{-line}:\text{even charged channels}.
\end{equation}
The global class is the channel supremum.  The relevant subregular orbit is
\begin{equation}
\label{eq:bp-ledger-subregular-ds}
\mathcal W_n^{(2)}(k)
=H^0_{\mathrm{DS}}\!\left(V_k(\mathfrak{sl}_n),f_{\mathrm{sub}}\right),
\qquad
f_{\mathrm{sub}}\leftrightarrow (n-1,1),
\end{equation}
and $\mathcal W_3^{(2)}(k)=\mathrm{BP}_k$.

For this subregular line, Theorem~\ref{thm:pbw-slodowy-collapse} computes the
algebraic PBW bar input, and
Theorem~\ref{thm:subregular-appell-formula} supplies the Appell formula for the
singular $E$--$F$ coefficients.  The exact OPE-degree window is
\begin{equation}
\label{eq:bp-ledger-degree-window}
d_{\mathrm{OPE}}(\mathcal W_n^{(2)})\ge r
\qquad (2\le r\le n-1)
\end{equation}
on the centred symbolic normal form.  A statement about completed or minimal
bar-dual Taylor coefficients requires $H_{\mathrm{OPE/bar}}$.
\end{proposition}

\begin{remark}[Genus-$1$ curvature calculation required]
\label{rem:bp-ghost-subtraction-failure}
\index{Bershadsky--Polyakov!ghost subtraction}
The BP genus-$1$ curvature receives contributions from the charged ghost
system, the neutral fields, the improvement of the stress tensor, and mixed
collision channels.  A valid computation starts from the non-separating
genus-$1$ contraction in the BRST model and descends it through a specified
DS/bar comparison.  The central charge and the strong-generator spectrum
supply consistency checks; the non-separating contraction determines the
answer.
\end{remark}

\begin{theorem}[Bershadsky--Polyakov same-family companion;
\ClaimStatusConditional]
\label{thm:bp-koszul-self-dual}
\index{Bershadsky--Polyakov!same-family companion|textbf}
\index{non-principal W-algebra!conditional companion}
\index{hook-type duality!self-transpose}
\textup{[}Type signature:
$(\mathrm{Open}\cap\mathrm{Cat},\mathrm{factorisation},
\mathrm{level}=1\leftrightarrow 2,
\mathrm{hypothesis\ package}=k\ne -3;
H_{\mathrm{BP}}^{\mathrm{DS/bar}}=
\{\mathrm{subregular\ hook\ DS/bar\ transport};
\mathrm{finite\text{-}type\ bar\ dual\ or\ completed\ continuous\ nondegenerate\ Verdier\ pairing};
\mathrm{comparison\ at\ regular\ level}\})$\textup{]}
For $k\ne-3$, assume $H_{\mathrm{BP}}^{\mathrm{DS/bar}}$.  Then the
candidate Verdier--Koszul companion of
$\mathrm{BP}_k=\mathcal W^k(\mathfrak{sl}_3,f_{(2,1)})$ is the same
subregular family at reflected level:
\[
(\mathrm{BP}_k)^!\simeq\mathrm{BP}_{-k-6}.
\]
The self-transpose identity $(2,1)^t=(2,1)$ selects the family on the
right-hand side.  KSDual membership additionally requires an object-level
fixed-point equivalence with an involutive nondegenerate Verdier pairing and
remains open.

On the unconditional scalar lane one has
\[
c_{\mathrm{BP}}(k)=-\frac{(2k+3)(3k+1)}{k+3},
\qquad
c_{\mathrm{BP}}(k)+c_{\mathrm{BP}}(-k-6)=50.
\]
The modular quantities $\kappa_{\mathrm{BP}}(k)$ and
$K^\kappa_{\mathrm{BP}}$ await the genus-$1$ computation described in
Remark~\ref{rem:bp-ghost-subtraction-failure}.
\end{theorem}

\begin{proof}
The package $H_{\mathrm{BP}}^{\mathrm{DS/bar}}$ is exactly the premise
needed to apply Theorem~\ref{thm:hook-transport-corridor} to the
self-transpose partition $(2,1)$.  The scalar identity is
Proposition~\ref{prop:bp-complementarity-constant}.
\end{proof}

\begin{remark}[Exact symbolic identities]
\label{rem:bp-self-dual-verification}
The engine \texttt{compute/lib/bp\_koszul\_conductor\_engine.py}
evaluates the standard FKR identity $K^c_{\mathrm{BP}}=50$, the shifted
secondary identity $K^{c,\mathrm{shift}}_{\mathrm{BP}}=196$, the even parity
of all four generators, and the open status of the genus-$1$ invariants.
These exact symbolic evaluations establish the stated scalar formulas.
The conditional DS/bar equivalence uses the categorical package in the
theorem.
\end{remark}

\subsection{\texorpdfstring{$\mathcal W_4^{(2)}$}{W4(2)} and the first cubic obstruction}

\begin{proposition}[\texorpdfstring{$\mathcal W_4^{(2)}$}{W4(2)}
OPE in Feigin--Semikhatov normal form; \ClaimStatusProvedElsewhere]
\label{prop:w4-fs-normal-form-ope}
For $n=4$, Feigin--Semikhatov write
\cite[(A.4.2)]{FeiginSemikhatov} (the cubic $H^3$ term below is genuine
in the canonical normal form; minimal-model survival of canonical
degree~$3$ is conditional on the Prochazka--Rapcak finite-spin
$W_n^{(2)}$ closure and the Yamada finite presentation).
\begin{align}
E(z)F(w)
&=
\frac{(k+2)(2k+5)(3k+8)}{(z-w)^4}
+
\frac{4(k+2)(2k+5)H(w)}{(z-w)^3}
\notag\\
&\quad
+\frac{(k+2)\bigl(-(k+4)T(w)+6HH(w)+2(2k+5)\partial H(w)\bigr)}{(z-w)^2}
\notag\\
&\quad
+\frac{(k+2)\Bigl(
W(w)-\frac12(k+4)\partial T(w)
-\frac{4(k+4)}{3k+8}T H(w)
-\frac{8(11k+32)}{3(3k+8)^2}H^3(w)
\Bigr)}{z-w}
\notag\\
&\quad
+\frac{(k+2)\Bigl(
6\,\partial H\,H(w)
+\frac{4(26+17k+3k^2)}{3(3k+8)}\partial^2 H(w)
\Bigr)}{z-w}.
\label{eq:w4-ef}
\end{align}
The simple pole contains an explicit cubic term in $H$.
\end{proposition}

\begin{theorem}[Cubic OPE degree in
\texorpdfstring{$\mathcal W_4^{(2)}$}{W4(2)}; \ClaimStatusProvedHere]
\label{thm:w4-cubic}
\textup{[}Type signature:
$(\mathrm{Open}\cap\mathrm{Chain},\mathrm{chain\text{-}level},
\mathrm{level}=1\leftrightarrow 2,
\mathrm{hypothesis\ package}=(k+2)(11k+32)\ne 0;
\mathrm{Feigin\text{--}Semikhatov\ OPE\ normal\ form})$\textup{]}
For every $k$ with $(k+2)(11k+32)\neq0$ and $3k+8\neq0$,
\[
d_{\mathrm{OPE}}(\mathcal W_4^{(2)})\ge3.
\]
The witness is the $H^3$ coefficient in \eqref{eq:w4-ef}.
\end{theorem}

\begin{proof}
The coefficient of $H^3$ in \eqref{eq:w4-ef} is
\[
-\frac{8(k+2)(11k+32)}{3(3k+8)^2},
\]
which is nonzero under the stated genericity condition. This is a generator
degree-$3$ term in the fixed Feigin--Semikhatov normal form.
\end{proof}

\begin{remark}[Meaning of the cubic term]
The neutral weight-$3$ generator $W$ is characterized in the standard
normalization by being Virasoro primary and regular with respect to $H$. Any
redefinition preserving that normalization can alter linear and derivative
terms.  The pure $H^3$ coefficient is the invariant cubic witness in the
$H$-regular gauge. This is the first subregular OPE in
which the centred normal form contains a cubic generator monomial.  Its
interpretation as a nonzero transferred $m_3$ is conditional on
$H_{\mathrm{OPE/bar}}$.
\end{remark}

\section{Explicit \texorpdfstring{$W_5^{(2)}$}{W5(2)} and \texorpdfstring{$W_6^{(2)}$}{W6(2)} symbols}

The Appell formula gives explicit symbolic coefficient packets.

\subsection{\texorpdfstring{$\mathcal W_5^{(2)}$}{W5(2)}}

For $n=5$,
\[
h_5=q+\frac{4a_1+3a_2+2a_3+a_4}{5},
\]
and
\[
(v_0,\dots,v_4)
=
\frac{1}{5}
\Bigl(
-4a_1-3a_2-2a_3-a_4,\;
a_1-3a_2-2a_3-a_4,\;
a_1+2a_2-2a_3-a_4,\;
a_1+2a_2+3a_3-a_4,\;
a_1+2a_2+3a_3+4a_4
\Bigr).
\]
The normalized symbolic generators are
\[
\mathcal W^{\mathrm{ess}}_{5,3}:=e_3(v_0,\dots,v_4),
\qquad
\mathcal W^{\mathrm{ess}}_{5,4}:=e_4(v_0,\dots,v_4).
\]
The raw weight-$3$ and weight-$4$ symbols are
\begin{align}
C_{5,3}
&=
\mathcal W^{\mathrm{ess}}_{5,3}
+ 3\,e_2(v)\,h_5
+ 10\,h_5^3,
\label{eq:w5-c3}
\\[4pt]
C_{5,4}
&=
\mathcal W^{\mathrm{ess}}_{5,4}
+ 2\,\mathcal W^{\mathrm{ess}}_{5,3}\,h_5
+ 3\,e_2(v)\,h_5^2
+ 5\,h_5^4.
\label{eq:w5-c4}
\end{align}

\subsection{\texorpdfstring{$\mathcal W_6^{(2)}$}{W6(2)}}

For $n=6$,
\[
h_6=q+\frac{5a_1+4a_2+3a_3+2a_4+a_5}{6},
\]
and
\[
(v_0,\dots,v_5)
=
\frac{1}{6}
\Bigl(
-5a_1-4a_2-3a_3-2a_4-a_5,\;
a_1-4a_2-3a_3-2a_4-a_5,\;
a_1+2a_2-3a_3-2a_4-a_5,\;
a_1+2a_2+3a_3-2a_4-a_5,\;
a_1+2a_2+3a_3+4a_4-a_5,\;
a_1+2a_2+3a_3+4a_4+5a_5
\Bigr).
\]
The normalized symbolic generators are
\[
\mathcal W^{\mathrm{ess}}_{6,3}:=e_3(v_0,\dots,v_5),
\qquad
\mathcal W^{\mathrm{ess}}_{6,4}:=e_4(v_0,\dots,v_5),
\qquad
\mathcal W^{\mathrm{ess}}_{6,5}:=e_5(v_0,\dots,v_5).
\]
The raw symbolic coefficient packets are
\begin{align}
C_{6,3}
&=
\mathcal W^{\mathrm{ess}}_{6,3}
+4\,e_2(v)\,h_6
+20\,h_6^3,
\label{eq:w6-c3}
\\[4pt]
C_{6,4}
&=
\mathcal W^{\mathrm{ess}}_{6,4}
+3\,\mathcal W^{\mathrm{ess}}_{6,3}\,h_6
+6\,e_2(v)\,h_6^2
+15\,h_6^4,
\label{eq:w6-c4}
\\[4pt]
C_{6,5}
&=
\mathcal W^{\mathrm{ess}}_{6,5}
+2\,\mathcal W^{\mathrm{ess}}_{6,4}\,h_6
+3\,\mathcal W^{\mathrm{ess}}_{6,3}\,h_6^2
+4\,e_2(v)\,h_6^3
+6\,h_6^5.
\label{eq:w6-c5}
\end{align}

\begin{theorem}[Unbounded symbolic OPE degree in the subregular line;
\ClaimStatusProvedHere]
\label{thm:unbounded-canonical-degree}
\textup{[}Type signature:
$(\mathrm{Open}\cap\mathrm{Chain},\mathrm{chain\text{-}level},
\mathrm{level}=1\leftrightarrow 2,
\mathrm{hypothesis\ package}=\mathrm{generic\ level\ }k;
\mathrm{Feigin\text{--}Semikhatov\ centered\ Appell\ symbols};
\mathrm{Bell\text{-}recursion\ for\ raw\ singular\ coefficients})$\textup{]}
For every $n\ge3$, the centred Appell symbol of the weight-$r$ raw
$E$--$F$ coefficient contains a nonzero degree-$r$ monomial for
$2\le r\le n-1$.  Hence the symbolic OPE generator degree is unbounded as
$n\to\infty$.
\end{theorem}

\begin{proof}
For general $n$, Theorem~\ref{thm:subregular-appell-formula} shows that the
weight-$r$ raw coefficient has symbolic expansion
\[
C_{n,r}(h_n,v)=e_r(v)+\binom{n-r+1}{1}e_{r-1}(v)h_n+\cdots+\binom{n}{r}h_n^r.
\]
The term $\binom{n}{r}h_n^r$ is nonzero for every $2\leq r\leq n-1$,
which proves the symbolic statement.
\end{proof}

\begin{conjecture}[OPE/bar arity comparison; \ClaimStatusConjectured]
\label{conj:minimal-degree-subregular}
For generic~$k$, there is a package $H_{\mathrm{OPE/bar}}$ identifying the
centred OPE generator filtration with the arity filtration of a specified
completed bar-dual model of $\mathcal W_n^{(2)}(k)$.  Under that comparison,
the degree-$r$ Appell witnesses survive transfer for $2\le r\le n-1$.
\end{conjecture}

\section{Primary renormalization and the full VOA}

Theorem~\ref{thm:full-raw-coefficient-packet} gives the complete raw singular
coefficients in the $E_n$--$F_n$ OPE. To pass from these raw coefficients to
the standard strong generators, one must project each weight-$r$ coefficient to
an $H$-regular Virasoro primary plus a uniquely determined descendant packet.

\begin{definition}[Orthogonal stress tensor]
\label{def:orthogonal-stress}
Write
\[
\ell_n(k)=\frac{n-1}{n}k+n-2,
\qquad
T_n^\perp:=T_n-\frac{1}{2\ell_n(k)}H_nH_n.
\]
Then $T_n^\perp$ has regular OPE with $H_n$.
\end{definition}

\begin{theorem}[Triangular primary-renormalization theorem; \ClaimStatusProvedHere]
\label{thm:triangular-primary-renormalization}
Fix $r\geq 2$. Let $\mathcal N_r$ be the vector space of $H$-neutral
normally ordered differential monomials of total conformal weight $r$ built from
\[
H_n,\quad T_n^\perp,\quad W_{n,3},\dots,W_{n,r-1},
\]
and derivatives. Let $\mathcal D_r\subset \mathcal N_r$ be the descendant
subspace generated by derivatives of lower-weight neutral fields and by
normally ordered products of lower-weight primaries with $H_n$ and its
derivatives. Let $\mathcal P_r\subset \mathcal N_r$ be the subspace of
$H$-regular Virasoro primaries of weight $r$.

Then, away from the determinant locus of the Shapovalov--Gram matrix of the
neutral weight-$r$ sector, there is a direct-sum decomposition
\[
\mathcal N_r=\mathcal P_r\oplus \mathcal D_r.
\]
Hence every raw coefficient $U_{n,r}^{\mathrm{raw}}$ admits a unique
decomposition
\[
U_{n,r}^{\mathrm{raw}}=W_{n,r}+D_{n,r},
\qquad
W_{n,r}\in \mathcal P_r,
\quad
D_{n,r}\in \mathcal D_r.
\]
The coefficients of $D_{n,r}$ are rational functions of $k$, obtained by
solving a triangular linear system against the neutral Shapovalov--Gram matrix.
\end{theorem}

\begin{proof}
Order a basis of $\mathcal N_r$ by derivative degree and then by the total
number of non-$H$ primaries. The Virasoro-primary and $H$-regularity
conditions are linear in this basis. Because derivatives raise derivative
degree, the action of $L_1$ and of the positive Heisenberg modes is triangular
with respect to the chosen ordering. For generic $k$ the corresponding
Shapovalov--Gram matrix is invertible on the descendant sector, hence the kernel
of the positive modes splits as a complement to the descendant subspace:
\[
\mathcal N_r=\mathcal P_r\oplus\mathcal D_r.
\]
Projecting $U_{n,r}^{\mathrm{raw}}$ to $\mathcal P_r$ gives the unique primary
$W_{n,r}$; the difference is $D_{n,r}\in\mathcal D_r$. All coefficients are
therefore rational in $k$.
\end{proof}

\begin{remark}[What ``every coefficient'' means]
Theorem~\ref{thm:full-raw-coefficient-packet} writes every singular coefficient
in closed recursive form. Theorem~\ref{thm:triangular-primary-renormalization}
then converts each raw coefficient into the standard primary generator plus its
entire lower-derivative packet. This is the conceptual content of ``writing
all lower-derivative, $k$-dependent, primary-renormalized coefficients'':
the recursion is exact, and the projection is an explicit triangular linear
algebra problem in each weight.
\end{remark}

\subsection{Weights two and three in closed form}

For $r=2$ the projection is immediate:
\[
U_{n,2}^{\mathrm{raw}}
=
\lambda_{n-3}(n,k)
\left(
\frac{n(n-1)}{2}H_nH_n
+
\frac{n((n-2)(k+n-1)-1)}{2}\partial H_n
-
(k+n)T_n
\right),
\]
so the $H$-regular part is precisely $T_n^\perp$.

For $r=3$, Feigin--Semikhatov's explicit formula gives
\begin{align}
U_{n,3}^{\mathrm{raw}}
&=
\lambda_{n-3}(n,k)
\Bigl(
W_{n,3}
-(k+n)\Bigl(\frac12\partial T_n^\perp+\frac{1}{\ell_n(k)}H_nT_n^\perp\Bigr)
\Bigr)
\notag\\
&\quad
+
\lambda_{n-2}(n,k)
\Bigl(
\frac{n}{6\ell_n(k)^2}H_n^3
+
\frac{n}{2\ell_n(k)}\partial H_n\,H_n
+
\frac{n}{6}\partial^2H_n
\Bigr).
\label{eq:weight-three-renormalized}
\end{align}
Thus the complete lower-derivative packet at weight $3$ is already universal.

\subsection{Weights four and five: exact descendant bases}

At weight $4$, a natural spanning set for the descendant packet is
\begin{align*}
\mathbb B_4
=
\Bigl\{
&\partial W_{n,3},\;
H_nW_{n,3},\;
\partial^2T_n^\perp,\;
H_n\partial T_n^\perp,\;
(\partial H_n)T_n^\perp,\;
H_n^2T_n^\perp,\;
T_n^\perp T_n^\perp,
\\
&\partial^3H_n,\;
H_n\partial^2H_n,\;
(\partial H_n)^2,\;
H_n^2\partial H_n,\;
H_n^4
\Bigr\}.
\end{align*}
Therefore the primary-renormalized weight-$4$ coefficient has the form
\[
U_{n,4}^{\mathrm{raw}}=W_{n,4}+\sum_{b\in\mathbb B_4}\alpha_b(n,k)\,b,
\]
with uniquely determined rational functions $\alpha_b(n,k)$.

At weight $5$, one similarly obtains a finite descendant basis generated by
\[
\partial^2W_{n,3},\;
H_n\partial W_{n,3},\;
(\partial H_n)W_{n,3},\;
H_n^2W_{n,3},\;
\partial W_{n,4},\;
H_nW_{n,4},
\]
together with all $H_n$-neutral composites built from $T_n^\perp$ and the
Heisenberg descendants of total weight $5$. Hence
\[
U_{n,5}^{\mathrm{raw}}=W_{n,5}+\sum_{b\in\mathbb B_5}\beta_b(n,k)\,b,
\]
again with uniquely determined rational functions $\beta_b(n,k)$.

\begin{remark}[Explicit expansion versus exact determination]
The complete rational expressions for the coefficients
$\alpha_b(n,k)$ and $\beta_b(n,k)$ are lengthy. For repository purposes the
proved datum is the exact Bell recursion together with the triangular
projection mechanism. This is stronger than a page-filling coefficient dump:
it is a first-principles construction of the full coefficient packets.
\end{remark}

\section{Classification package beyond the principal case}

The preceding sections divide the non-principal problem into two layers.

\begin{enumerate}[label=(\alph*)]
\item \emph{Filtered bar input.} The PBW--Slodowy theorem computes the
 associated-graded bar page.  Completed Koszulity is the conditional output of
 $H_{\mathrm{PBW}}^{\mathrm{bar}}$.
\item \emph{OPE nonlinearity.} The centred normal form has the exact pattern
 \[
 d_{\mathrm{OPE}}(\mathcal W_3^{(2)})=2,\qquad
 d_{\mathrm{OPE}}(\mathcal W_4^{(2)})\ge3,\qquad
 \sup_n d_{\mathrm{OPE}}(\mathcal W_n^{(2)})=\infty.
 \]
 A bar-dual arity statement is the comparison conjecture
 \ref{conj:minimal-degree-subregular}.
\end{enumerate}

Thus the hook/subregular family forms a coherent staircase of OPE
nonlinearity.  The derived meaning of that staircase is a sharply stated open
problem.

\begin{conjecture}[Hook-type canonical-versus-minimal coincidence; \ClaimStatusConjectured]
\label{conj:hook-canonical-minimal}
For generic level, $H_{\mathrm{OPE/bar}}$ extends functorially along the
hook-reduction graph and identifies the OPE generator filtration with the
arity filtration of a minimal transferred completed dual.
\end{conjecture}

\begin{remark}[Geometric meaning]
The PBW side of the story is controlled by Slodowy slices and their arc spaces.
The reduction-to-hook-type side is controlled by inverse Hamiltonian reduction
charts and reduction-by-stages along the nilpotent-orbit order. The proposed
homological interpretation asks whether the higher Taylor coefficients of a
specified completed dual are detected by the nonlinear normal-ordering terms
after transport through $H_{\mathrm{OPE/bar}}$.
\end{remark}

\section{Outlook}

Three concrete problems remain.

\begin{enumerate}[label=(\roman*)]
\item Carry the triangular projection of
Theorem~\ref{thm:triangular-primary-renormalization} out explicitly in weights
$4$ and $5$, thereby producing the full primary-renormalized coefficient packets
for $\mathcal W_5^{(2)}$ and $\mathcal W_6^{(2)}$ as concrete rational
functions of $k$.
\item Construct the OPE/bar comparison by proving
Conjecture~\ref{conj:minimal-degree-subregular} or, more ambitiously,
Conjecture~\ref{conj:hook-canonical-minimal}.
\item Transport the same machinery through the type-$A$ inverse-reduction
network and verify the DS/bar edge package.
\end{enumerate}

The exact output of the chapter is the PBW associated-graded page, the full
Bell recursion for the raw singular packet, the Miura--Appell symbols, and the
standard BP central-charge conductor.  The remaining obligations are the
completed bar comparison, the genus-$1$ BP curvature, and DS/bar transport.

\section{Nilpotent transport and the even-nilpotent dichotomy}

The reduction graph underlying the hook-type corridor
(Theorem~\ref{thm:hook-transport-corridor}) provides the irreducible spine
of the type-$A$ reduction graph. The next question is \emph{coverage}:
for which ranks does the transport-closure of the hook vertices exhaust all
partitions?

\begin{proposition}[Finite-rank transport audit for type $A$;
\ClaimStatusConditional]
\label{prop:nilpotent-transport-typeA}
\index{nilpotent transport!type A}
\index{reduction graph!transport closure}
\textup{[}Type signature:
$(\mathrm{Open}\cap\mathrm{Cat},\mathrm{factorisation},
\mathrm{level}=1,
\mathrm{hypothesis\ package}=\mathfrak{sl}_N\mathrm{\ for\ }2\le N\le 8;
\mathrm{generic\ level};
\mathrm{Fehily\ \textup{(}hook\textup{)} +\ Creutzig\text{--}Linshaw\text{--}Nakatsuka\text{--}Sato\ \textup{(CLNS)}\ +\ }
\mathrm{Creutzig\text{--}Fasquel\text{--}Linshaw\text{--}Nakatsuka\ \textup{(CFLN)}};
\mathrm{Genra\text{--}Juillard\ stages};
\mathrm{Butson\text{--}Nair\ inverse\ Hamiltonian\ reduction})$\textup{]}
For $N=2,\ldots,8$, a finite combinatorial calculation shows that every
partition of~$N$ lies in the closure-ordering component of a hook partition.
Promotion of that Hasse-graph statement to the reduction graph~$\Gamma_N$
requires a proved reduction or inverse-reduction functor for every chosen edge.
Promotion further to Koszul transport requires
$H_{\mathrm{hook}}^{\mathrm{DS/bar}}$ on every edge.
\end{proposition}

\begin{proof}
The closure-ordering calculation is finite and direct.  The cited works
construct substantial families of reduction and inverse-reduction edges,
including the hook spine.  The statement that those results realise every
covering edge through rank~$8$, and that each realised edge commutes with the
completed chiral bar construction, is the conditional part of the
proposition.
\end{proof}

\begin{remark}[Even and self-transpose orbits as audit surfaces]
\label{rem:even-nilpotent-dichotomy}
\index{nilpotent orbit!even vs non-even}
\index{complementarity!even-nilpotent dichotomy}
Even and self-transpose partitions give efficient tests of the KRW central
charge under orbit transpose and level reflection.  These tests determine
candidate scalar conductors in a fixed conformal-vector convention.  A
classification of all level-independent sums requires a uniform evaluation of
the KRW formula across good gradings.  A modular complementarity theorem
requires, in addition, the genus-$1$ curvature and DS/bar/Verdier transport.
For the self-transpose BP partition $(2,1)$, the standard scalar audit gives
$K^c_{\mathrm{BP}}=50$.
\end{remark}

\begin{remark}[Beyond hook type: the rectangular partition $(2,2)$ in $\mathfrak{sl}_4$]
\label{rem:beyond-hook-rectangular}
\index{W-algebra@$\mathcal{W}$-algebra!beyond hook type}
\index{W-algebra@$\mathcal{W}$-algebra!(2,2) rectangular}
\index{DS reduction!beyond hook type}
The hook-type corridor
(Theorem~\ref{thm:hook-transport-corridor}) covers all
partitions of the form $(N{-}r, 1^r)$.
The first partition lying outside this family is $(2,2)$ in
$\mathfrak{sl}_4$: a self-transpose rectangular even nilpotent.
The following data organise the first audit beyond the hook spine for
$\mathcal{W}^k(\mathfrak{sl}_4, f_{(2,2)})$
at generic level~$k \neq -4$:

\begin{enumerate}[label=\textup{(\roman*)},nosep]
\item \emph{Filtered surface.}
 The Kazhdan--Li theorem supplies the Slodowy associated graded.  The
 PBW--Slodowy bar page follows from
 Theorem~\ref{thm:pbw-slodowy-collapse}.  Completed chiral Koszulity is
 conditional on $H_{\mathrm{PBW}}^{\mathrm{bar}}$.

\item \emph{Central charge and generator spectrum.}
 In the displayed good-grading convention, direct evaluation of the KRW
 formula gives
 \[
 c(k)=\frac{15k}{k+4}-12k-8
 =\frac{-12k^2-41k-32}{k+4},
 \]
 The strong generating set has $7$ even generators
 ($3$ at weight~$1$, $4$ at weight~$2$).  The modular characteristic requires a separate genus-$1$
 curvature calculation.

\item \emph{Candidate companion.}
 Since $(2,2)$ is self-transpose, orbit transpose selects the same
 $W$-family at candidate reflected level $k^\vee=-k-8$.  The
 rectangular DS--bar/Verdier exchange package required to construct this
 companion is open.  KSDual membership is a further fixed-point statement.

\item \emph{BRST complex.}
 The nilradical $\mathfrak{n}_+$ under the $(2,2)$-grading
 has $\dim(\mathfrak{n}_+) = 4$, all roots at
 $\operatorname{ad}(x)$-grade~$1$, and
 $\mathfrak{n}_+$ is abelian. This is in contrast to generic
 non-hook partitions (e.g.\ $(3,2)$ in~$\mathfrak{sl}_5$),
 where $\mathfrak{n}_+$ is non-abelian with half-integer
 grades.

\item \emph{Shadow audit.}
 The weight-$1$ and weight-$2$ generator spectrum gives a finite OPE-channel
 diagnostic.  The global shadow depth is determined by the modular MC
 coefficients and therefore requires their direct calculation.
\end{enumerate}

\noindent\textit{Audit classes beyond hooks.}
The following partition types organise the comparison problem:
\begin{center}
\renewcommand{\arraystretch}{1.15}
\begin{tabular}{llll}
\toprule
\textbf{Class} & \textbf{Condition} &
 \textbf{$c + c'$ status} & \textbf{DS--KD status} \\
\midrule
$\mathsf{A}$ (hook) & $\lambda = (N{-}r, 1^r)$ &
 family-by-family & conditional on $H_{\mathrm{hook}}^{\mathrm{DS/bar}}$ \\
$\mathsf{B}$ (self-tr.\ rect.) & $\lambda = \lambda^t$, rectangular &
 KRW audit & open comparison \\
$\mathsf{C}$ (general self-tr.) & $\lambda = \lambda^t$, nonrectangular &
 KRW audit & open comparison \\
$\mathsf{D}$ (general pair) & $\lambda \neq \lambda^t$ &
 KRW audit & open comparison \\
\bottomrule
\end{tabular}
\end{center}

\noindent
The $(2,2)$ partition is the first rectangular audit surface;
$(3,2)$ in $\mathfrak{sl}_5$ is the first explicit non-hook surface with a
non-abelian reduction algebra.  In each class the associated-graded Slodowy
description, completed bar collapse, modular curvature, and DS--KD comparison
remain distinct obligations.
\end{remark}

\begin{remark}[Scope of the subregular and minimal routes]
\label{rem:subregular-minimal-scope}
\index{W-algebra@$\mathcal{W}$-algebra!subregular route!scope}
\index{W-algebra@$\mathcal{W}$-algebra!minimal route!scope}
Two informal decomposition routes often attached to non-principal
$\mathcal{W}$-algebras deserve explicit scoping.
\begin{enumerate}[label=\textup{(\roman*)}]
\item \emph{Subregular via a principal--triplet--ghost decomposition.}
The Feigin--Semikhatov free-field and screening realisations provide exact
presentations of the subregular family in the ranges established in the cited
literature.  A tensor-product decomposition into a smaller principal
$W$-algebra, a current algebra, and a free field is a family-specific theorem
and must be stated with its level map and screening cohomology.  This chapter
uses only the BP presentation and the general Bell/Miura formulas proved above.
\item \emph{Minimal via coset.}
In type~$A$, the minimal and subregular orbits coincide for
$\mathfrak{sl}_3$ and separate from rank~$4$ onward.  Coset descriptions are
therefore indexed by the orbit, the level transformation, and the precise
universal or simple quotient.  Completed Koszulity and DS/bar commutation are
the subsequent comparison problems.
\item \emph{Screening-kernel Koszul-locus preservation.}
A screening presentation identifies a vertex subalgebra or a cohomology object
inside a free-field parent.  Chiral Koszulity passes through that presentation
only after exactness of the screening complex and compatibility of its
filtration with the chiral bar differential are proved.  These clauses are the
screening part of $H_{\mathrm{hook}}^{\mathrm{DS/bar}}$.
\end{enumerate}
\end{remark}

\section{Boundary of the chapter}
\label{sec:subreg-supplement}

\begin{remark}[Separation from the K3 moonshine programme]
\label{rem:subreg-anomalous-classes}
The present chapter concerns affine $W$-algebras obtained by quantum
Drinfeld--Sokolov reduction and their prospective chiral bar companions.  The
K3, Mathieu-moonshine, and paramodular denominator constructions occupy a
separate CY/BKM row with their own categorical source and comparison functors.
Any future bridge between the two programmes must name a functor from the
subregular factorization category to the K3 Hall--Drinfeld datum and verify its
effect on OPE, trace, and Verdier pairing.
\end{remark}
